%! Setting Up a Test for a Population Mean — Unit 7, Lesson 7.4 — Warm-Up (Answer Key)
\documentclass[10pt]{article}
\usepackage{apstats-article}
\usepackage{apstats-key}   % pulls in -boxes; do NOT also load -boxes
\newcommand{\TallMath}[1]{$\displaystyle #1\rule[-1.4em]{0pt}{3.2em}$}

\begin{document}

\pageheader{Unit 7, Lesson 7.4}{Warm-Up --- Answer Key}
\namedateperiod

\vspace{0.15in}

\begin{notesbox}{Spiral Review --- The $t$-Distribution \& Sampling Distributions (Unit 7, Lessons 7.1--7.3)}

\textbf{1.} A $t$-distribution with 14 degrees of freedom is used to model the sampling
distribution of $\bar x$ from a sample of size $n = 15$.

\begin{enumerate}[label=\alph*.]
  \item How does the shape of a $t$-distribution compare to the standard normal ($z$) distribution?
        Circle the best description.

        \quad (A) Shorter tails, taller peak \quad (B) Identical \quad \ans{(C) Heavier tails, shorter peak} \quad (D) Skewed right

  \item As degrees of freedom increase, the $t$-distribution becomes \ans{closer to / approaches}
        the standard normal distribution.

  \item Why do we use a $t$-distribution instead of a $z$-distribution when making inferences
        about a population mean? \ansline{Because $\sigma$ is unknown; we estimate it with $s$, which adds extra variability --- the $t$-distribution accounts for this uncertainty.}
\end{enumerate}

\vspace{0.12in}
\textbf{2.} A school nutritionist wants to test whether the mean number of grams of sugar in
students' lunches differs from the district target of 30 g.

\begin{enumerate}[label=\alph*.]
  \item She knows $\bar x = 33.4$ g and $s = 5.2$ g from a random sample of $n = 22$ lunches,
        but $\sigma$ is unknown. Should she use a $z$-test or a $t$-test? \ans{$t$-test}

        Explain: \ansline{$\sigma$ is unknown; we only have $s$, so a $t$-test is required.}

  \item State the degrees of freedom she would use. \quad $df = n - 1 = $ \ans{21}

  \item The sample is taken from a population of 640 students.
        Is the 10\% condition met? \ans{Yes} Show: \ans{$22 \le 0.10 \times 640 = 64$}
\end{enumerate}

\end{notesbox}

\end{document}
